% ============================================================
% COMPUTATIONAL TOXICOLOGY – Special Issue: "Advancing FAIR and
% Integrative PBK Modelling for Next-Generation Risk Assessment"
% Submission deadline: 31 July 2026
% Article type to select in Editorial Manager: 'VSI: FAIR and PBK'
%
% FRAMING NOTES FOR AUTHORS:
% This paper (PBKFAIRModel + FAIRDatabase) is an excellent fit for Theme 1
% (FAIR Data and Model Standardization) and Theme 5 (Life-Stage and
% Population Variability) of the special issue. The PFAS early-life exposure
% model also directly addresses the special issue's explicit mention of PFAS as an emerging contaminant of interest.
%
% SUGGESTED FRAMING STRATEGY:
% - Lead with the FAIR-compliance angle: SBML encoding, open-source
%   web platform, role-based provenance tracking. This is the primary
%   contribution that distinguishes the paper.
% - Secondary angle: early-life (life-stage) pharmacokinetics for PFAS.
%   Explicitly position the breastfeeding/prenatal model as a demonstration of life-stage and population variability capabilities.
% - Tertiary angle: Integration with omics/microbiome data infrastructure
%   (FAIRDatabase), touching Theme 4 (New Approach Methodologies / multi-omics).
%
% MAIN MANUSCRIPT vs. SUPPLEMENTARY FILE SPLIT:
%
%   MAIN MANUSCRIPT should contain:
%     1. Abstract (≤250 words – currently too long, needs trimming)
%     2. Highlights (3–5 bullets, max 85 chars each)
%     3. Graphical Abstract (REQUIRED – currently missing, must be added)
%     4. Keywords (1–7, simple single-concept terms)
%     5. Introduction – motivation, FAIR gap, PFAS context, paper outline
%     6. Methodology:
%        - Brief model description (conceptual, not full ODEs)
%        - SBML encoding rationale (FAIRness argument)
%        - Integration architecture overview (FAIRDatabase)
%        - UI and API overview (summary table of endpoints)
%     7. Results & Discussion:
%        - Simulation outputs for 4 breastfeeding scenarios (fill in table!)
%        - FAIR compliance evaluation (checklist-style)
%        - Comparison with httk, IRRA frameworks
%        - Limitations
%     8. Conclusion
%     9. Declarations (CRediT, Funding, COI, AI use)
%     10. References (numbered, square brackets [1,2] style)
%
%   SUPPLEMENTARY FILE should contain:
%     S1. Full SBML module descriptions and ODE listing
%     S2. Complete table of all 64 state variables
%     S3. Complete API endpoint table (all 11 endpoints)
%     S4. Parameter values for PFOA and PFOS
%     S5. Full breastfeeding scenario parameterization table
%     S6. Installation and dependency instructions
%     S7. Python runner code listing (simulate_scipy.py excerpt)
%     S8. Role-based access control detailed table
%     S9. Subsubsubsection-level UI details (layout, CSV export, etc.)
%
% NOTE: The current manuscript has too many low-level technical details
% (package structures, Flask blueprint registrations, Jinja2 template names, etc.) in the main body. These belong in the supplementary file to keep the main text focused on scientific contribution and FAIRness demonstration.
% ============================================================

\documentclass[preprint,12pt]{elsarticle}

% ---- Elsevier elsarticle packages ----
\usepackage{graphicx}
\usepackage{multirow}
\usepackage{amsmath,amssymb,amsfonts}
\usepackage{mathrsfs}
\usepackage[title]{appendix}
\usepackage{xcolor}
\usepackage{textcomp}
\usepackage{booktabs}
\usepackage{algorithm}
\usepackage{algorithmicx}
\usepackage{algpseudocode}
\usepackage{listings}
\usepackage{url}
\def\UrlFont{\rmfamily}
\usepackage{enumitem}
\usepackage{array}
\usepackage{lineno}

% Elsevier natbib for numbered [1] style citations
\usepackage[numbers,sort&compress]{natbib}

% ====================== JOURNAL METADATA ===================
% NOTE: Use elsarticle document class for Elsevier submission.
% Select journal: Computational Toxicology (ISSN 2468-1113)
% In Editorial Manager, select article type: 'VSI: FAIR and PBK'
% ============================================================

\journal{Computational Toxicology}

\begin{document}

\begin{frontmatter}

% ====================== TITLE ==============================
% NOTE: Title should be concise and informative. Avoid abbreviations
% unless firmly established (PFAS, SBML, FAIR are acceptable).
% Consider making the FAIR angle more prominent in the title for
% the special issue. Suggested alternative:
% "PBKFAIRModel: A FAIR-Compliant Web Platform for Early-Life
%  PFAS Exposure Assessment Using SBML-Encoded PBK Modelling"
% ============================================================
\title{\bfseries Integrating a PBK Module in FAIRDatabase Web Platform for Exposure Assessment: PFAS Case Study Across Pregnancy and Early Life}

% OLD TITLE: Integrating a FAIR-Compliant PBK Model into FAIRDatabase: Implementation and Web Interface for PFAS Exposure Assessment in Children

% ====================== AUTHORS ============================
% NOTE: Elsevier elsarticle uses \author and \address commands.
% Add ORCID IDs if available -- strongly encouraged by the journal.
% Mark corresponding author with \corref.
% ============================================================
\author[1,2]{Tu-Ky Ly}
\author[1,2,3,4]{Roland V. Bumbuc}
\author[2]{Senna Wiedeman}
\author[1,2]{Vivek M. Sheraton\corref{cor1}}
\ead{v.s.muniraj@uva.nl}
\cortext[cor1]{Corresponding author}

\address[1]{Computational Science Lab, Informatics Institute,
University of Amsterdam, Amsterdam, The Netherlands}
\address[2]{Informatics Institute, University of Amsterdam,
Amsterdam, The Netherlands}
\address[3]{Department of Plastic, Reconstructive and Hand Surgery,
Amsterdam Movement Sciences (AMS) Institute, Amsterdam UMC,
Location VUmc, Amsterdam, The Netherlands}
\address[4]{Department of Molecular Cell Biology and Immunology,
Amsterdam Infection and Immunity (AII) Institute, Amsterdam UMC,
Location VUmc, Amsterdam, The Netherlands}

\end{frontmatter}

% ====================== DATE ================================
% NOTE: Remove \date from the frontmatter; elsarticle handles dates.
% ============================================================

\begin{linenumbers}

% ====================== GRAPHICAL ABSTRACT =================
% NOTE: A GRAPHICAL ABSTRACT IS REQUIRED by Computational Toxicology.
% It is currently MISSING from the manuscript.
% Create a single-panel figure (~500x400 px, min 300 DPI) showing:
%   - A schematic of the PBKFAIRModel workflow:
%     [SBML model] --> [Python runner] --> [FAIRDatabase] --> [Web UI]
%   - Small inset: example venous PFAS concentration curve for the
%     4 breastfeeding scenarios
% Submit as a separate image file (TIFF or EPS preferred, JPEG acceptable).
% ============================================================

% ====================== HIGHLIGHTS =========================
% NOTE: Required. Max 5 bullets, max 85 characters each (including spaces).
% Current bullets are within acceptable length but verify character count.
% Highlights must appear BEFORE the abstract in the submission system,
% uploaded as a separate file or in the cover letter system.
% ============================================================
\section*{Highlights}
\begin{itemize}
  \item FAIR-compliant PBK model for PFAS encoded in SBML and deployed via Python.
  \item FAIRDatabase extended from omics to computational toxicokinetic modeling.
  \item Web UI enables scenario configuration and real-time result visualization.
  \item Four breastfeeding scenarios quantify early-life PFAS body burden.
  \item Role-based access control enables collaborative simulation workflows.
\end{itemize}

% ====================== ABSTRACT ===========================
% NOTE: Max 250 words (STRICT – current abstract is ~185 words, which is fine).
% No citations in abstract. No non-standard abbreviations without definition.
% Reframe opening sentence to mention FAIR principles first – this aligns
% with the special issue framing. End with a clear impact statement.
% Avoid mentioning "MCSim" without explanation; reviewers may not know it.
% ============================================================
\begin{abstract}
Physiologically-based kinetic (PBK) models are essential tools for
quantifying internal chemical doses in risk assessment, yet their
adoption is limited by siloed implementations and the absence of
standardized, web-accessible platforms. We present PBK FAIR (Findable, Accessible, Interoperable, Reusable) Model, a
fully FAIR-compliant 
PBK modelling framework integrated into FAIRDatabase for computational
toxicology. The lifetime multi-compartment PBK model simulates
per- and polyfluoroalkyl substance (PFAS) toxicokinetics from birth
through early childhood, explicitly capturing prenatal exposure via
placental transfer, postnatal exposure via breastfeeding, and dietary
intake. The original model was translated into Systems Biology Markup
Language (SBML) Level 3 Version 2, ensuring machine-readability and
interoperability with community-standard tools. A Python-based
simulation engine integrated as a Flask blueprint within FAIRDatabase
extends the platform from omics data management to toxicokinetic
modelling. The web interface allows users to configure exposure
scenarios, chemical parameters, and birth year with real-time
visualization of venous plasma concentration profiles. Four
breastfeeding scenarios (none, 6 months, 1 year, 3 years) demonstrate
the model's capacity to quantify early-life PFAS body burden
trajectories. This work advances FAIR computational toxicology by
providing an open-source, web-accessible PBK platform with persistent
storage and provenance tracking, directly supporting next-generation
risk assessment of PFAS in vulnerable populations.

\textbf{Keywords:} PBK model; FAIR principles; PFAS; early-life exposure;
SBML; computational toxicology; web platform; risk assessment
\end{abstract}

% ====================== KEYWORDS ===========================
% NOTE: Journal requires 1–7 keywords. Use single-concept terms.
% Avoid multi-word phrases with "and" or "of" where possible.
% Suggested keywords aligned with special issue keywords:
% PBK modelling; FAIR principles; PFAS; risk assessment; SBML;
% early-life exposure; web platform
% The current keyword list has 8 items and includes "pregnancy" and
% "microbiome" which are secondary to the main contribution -- consider
% replacing with "risk assessment" and "open science".
% ============================================================

% ====================== INTRODUCTION =======================
% NOTE: Introduction should be ~600–900 words (approx. 1.5 columns).
% Current introduction is well-structured. Strengthen by:
%  1. Adding an explicit paragraph on the OECD validation framework
%     for PBK models – this is directly relevant to the special issue
%     and Theme 3 (Regulatory Relevance).
%  2. Mention the European PARC initiative (aligns with special issue
%     "Alignment with Global Initiatives" section).
%  3. Reference Aude Ratier (guest editor of this special issue) more
%     prominently -- her 2024 PFAS lifetime PBK model paper is already
%     cited and is the direct basis of this work.
%  4. End introduction with a clear paragraph listing the paper's
%     contributions in the context of the special issue themes.
% ============================================================
\section{Introduction}\label{sec1}

The assessment of chemical exposure in human populations, particularly
during critical windows of development such as pregnancy and early
childhood, requires sophisticated computational approaches.
Physiologically-based kinetic (PBK) models have emerged as essential
tools for quantifying internal dose metrics, predicting toxicokinetic
behaviour, and informing risk assessment decisions~\cite{rabier2008,
welch2005, andersen2003}. However, the accessibility and reusability
of these models remain limited by siloed implementations, proprietary
formats, and the absence of standardized interfaces for integration
with data management platforms.

The FAIR (Findable, Accessible, Interoperable, Reusable) principles
\cite{wilkinson2016} have transformed data management across scientific
disciplines, yet their application to computational models in toxicology
remains incomplete. Regulatory bodies, including the OECD, have
published guidance documents for the validation of PBK models, yet a
persistent gap remains between model development and its practical
deployment in accessible, reproducible platforms that meet FAIR
standards. International initiatives such as the European Partnership
for the Assessment of Risks from Chemicals (PARC) and efforts to
develop PBK ontologies further underscore the urgency of
FAIR-compliant model infrastructure.

% NOTE: Add 1-2 paragraphs here on:
%  - PFAS as an environmental health concern; vulnerable populations
%    (pregnant women, infants) – keep from current draft, tighten
%  - Gap: no unified FAIR platform combining PBK execution + data storage

FAIRDatabase~\cite{dorst2024, roman2026} was originally developed as
an open-source, privacy-compliant infrastructure for human microbiome
data. In this work, we extend FAIRDatabase to computational
toxicokinetic modelling through PBKFAIRModel, a lifetime
multi-compartment PBK model for PFAS exposure encoded in SBML and
deployed via a web-based platform.

% NOTE: End introduction with a structured contribution paragraph, e.g.:
% "This paper makes the following contributions: (i) SBML encoding of
% the Ratier et al. (2024) lifetime PFAS PBK model; (ii) Python-native
% simulation engine; (iii) integration into FAIRDatabase; (iv)
% demonstration via four breastfeeding scenarios."

\section{Methodology}\label{sec2}

% ====================== METHODOLOGY ========================
% NOTE: This section needs significant restructuring for the journal.
% The current draft contains very deep technical detail (Flask blueprint
% registration, Jinja2 template names, Python package paths) that belongs
% in the Supplementary File (see S6, S7 below).
% Main text should present:
%   - Model conceptual description (compartments, ODEs, mass-balance)
%   - SBML encoding rationale (why SBML = FAIR)
%   - Integration architecture (conceptual, not code-level)
%   - UI overview (what a user can do, not how the HTML is structured)
% ============================================================

\subsection{Data Source and Model Selection}

% NOTE: Keep this subsection largely as-is. Add a sentence explicitly
% linking to OECD TG 417 or the OECD PBK model guidance document to
% strengthen the regulatory relevance framing.

The PBK model implemented in this study is a lifetime multi-compartment
model for PFAS exposure in humans, originally developed by Beaudouin
et al. (2010)~\cite{beaudouin2010} and subsequently refined by Ratier
et al. (2024)~\cite{ratier2024}. The model simulates PFAS
toxicokinetics from birth through early childhood, explicitly
representing prenatal exposure via placental transfer, postnatal
exposure via breastfeeding, and dietary intake. PFOA
(perfluorooctanoic acid) was used as the reference chemical, with a
plasma half-life of 2.5 years and a free plasma fraction of 0.02.
Parameters for other PFAS, including PFOS (plasma half-life 3.5 years),
can be substituted by modifying partition coefficients in the SBML
parameter block.

\subsection{PBK Model Implementation}

\subsubsection{SBML Encoding}

The model was encoded in Systems Biology Markup Language (SBML)
\cite{hucka2003}, Level 3 Version 2, as the file
\texttt{lifetime\_pbpk.xml}. SBML was selected as the target format
because it is a community-standard, machine-readable representation
of biochemical network models that directly supports the FAIR
principles of Interoperability and Reusability. The SBML file serves
as the committed source of truth; all modules in the \texttt{sbml/}
directory generate and validate this file via
\texttt{generate\_sbml.py}.

The encoded model comprises 64 state variables (species) organized
across two physiological sub-models: (i) a \textbf{maternal sub-model}
with 23 organ compartments including adipose, liver, kidney, lung, gut,
brain, breast, uterus, placenta, amniotic fluid, and blood pools; and
(ii) a \textbf{foetal sub-model} with 19 foetal organ compartments,
with foetal body weight following the Luecke (1994) polynomial growth
trajectory~\cite{luecke1994}. Full module descriptions and a complete
listing of all 64 state variables are provided in the Supplementary
File (Tables S1--S2).

% NOTE: Table 1 (SBML modules) should be moved to the Supplementary File
% as Table S2. In the main text, refer to it as "See Table S2 in
% Supplementary File." This dramatically lightens the main text while
% preserving all technical detail for reproducibility.

\subsubsection{Python Simulation Engine}

Numerical simulation is performed by \texttt{simulate\_scipy.py},
which parses \texttt{lifetime\_pbpk.xml} using \texttt{python-libsbml}
and integrates the 64-ODE system using the LSODA algorithm from
\texttt{scipy.integrate}. The compiled ODE functions are cached as a
module-level singleton after initial parse ($\approx$2~s), so that
all four reference scenarios complete in approximately 5 seconds on
standard hardware. The runner exposes a simple Python API (see
Supplementary File, Section S7 for code listing).

\subsubsection{Breastfeeding Exposure Scenarios}

% NOTE: Table 2 (scenarios) is appropriate for the main text -- keep it.
% However, the full parameterization (StopBreastmilk_total, C_milk_input,
% etc.) belongs in Supplementary Table S5.

Four pre-configured exposure scenarios are defined, corresponding to
different breastfeeding durations for a male child born in 2007:

\begin{table}[h]
\centering
\caption{Reference breastfeeding scenarios simulated by the PBK model.}
\label{tab:scenarios}
\begin{tabular}{lll}
\toprule
\textbf{Label} & \textbf{Breastfeeding duration} & \textbf{Description} \\
\midrule
\texttt{no\_bf}  & None     & No breastfeeding \\
\texttt{bf\_6mo} & 6 months & Exclusive breastfeeding for 6 months \\
\texttt{bf\_1yr} & 1 year   & Breastfeeding for 1 year \\
\texttt{bf\_3yr} & 3 years  & Extended breastfeeding for 3 years \\
\bottomrule
\end{tabular}
\end{table}

Full parameterization of each scenario is provided in Supplementary
Table~S5.

\subsection{Integration into FAIRDatabase}

% NOTE: This subsection should focus on the ARCHITECTURE and the FAIR
% compliance argument, not the Flask implementation details.
% Implementation details (package paths, blueprint registration,
% requirements.txt changes) move to Supplementary Section S6.

The PBKFAIRModel was integrated into FAIRDatabase as a modular
extension, following the platform's layered Flask architecture.
FAIRDatabase is an open-source, privacy-compliant infrastructure
built on a Supabase-managed PostgreSQL backend, originally developed
for human microbiome research~\cite{dorst2024, roman2026}. The
integration extends the platform from omics data management to
computational toxicokinetic modelling, providing:

\begin{itemize}
  \item A self-contained Python package (\texttt{PBKFAIRModel/})
        containing the SBML model and simulation runner.
  \item A Flask blueprint (\texttt{/model}) exposing REST API endpoints
        for scenario execution, parameter set management, and artifact
        storage (see Supplementary Table~S3 for full endpoint list).
  \item A web-based user interface at \texttt{/model/ui} extending
        the FAIRDatabase dashboard.
\end{itemize}

Role-based access control (RBAC) with three roles (Admin, Curator,
Accessor) ensures secure collaborative workflows. Full RBAC
specifications are provided in Supplementary Table~S8.

\subsection{User Interface}

% NOTE: Keep this section brief in the main text. Move subsubsection-
% level detail (layout, artifact upload file formats, CSV export
% mechanics, Chart.js configuration) to Supplementary Section S9.

The PBKFAIRModel user interface provides an accessible, interactive
experience for configuring and visualizing PBK simulations without
requiring command-line access or programming expertise. Users configure
simulations through four inputs: breastfeeding scenario, chemical
half-life (years), dietary intake rate (ng/kg/min), and birth year.
Results are displayed as an interactive time-concentration profile
(Chart.js) with four summary statistics (peak concentration, age at
peak, final concentration, final age) and a full CSV export. Artifact
storage supports images, videos, and mesh files with role-gated access.

% ====================== RESULTS AND DISCUSSION =============
\section{Results and Discussion}\label{sec3}

% NOTE: The Results & Discussion section currently has placeholder text
% ("to be populated"). This MUST be completed before submission.
% Populate Table 3 with actual simulation outputs.
% Add 1–2 figures:
%   Figure 1 (REQUIRED): Time-concentration profiles for all 4 scenarios
%     on a single plot (venous PFAS concentration vs. age in years).
%   Figure 2 (RECOMMENDED): Screenshot of the FAIRDatabase web UI
%     showing the simulation interface and result visualization.
% ============================================================

\subsection{Simulation Results}

% NOTE: Fill in Table 3 with actual values from model execution.
% If PFOA simulation values from Ratier et al. 2024 are available,
% use those as validation benchmarks and compare your outputs.

Table~\ref{tab:results_summary} summarizes the key toxicokinetic
outcomes for the four breastfeeding scenarios (male child, birth
year 2007, PFOA reference parameters).

\begin{table}[h]
\centering
\caption{Summary toxicokinetic results for four breastfeeding scenarios
(male, birth year 2007, PFOA: $t_{1/2}$=2.5 yr, free fraction=0.02).}
\label{tab:results_summary}
\begin{tabular}{lcccc}
\toprule
\textbf{Scenario} & \textbf{Peak C\textsubscript{ven} (mg/L)} &
\textbf{Age at peak (yr)} & \textbf{Final C\textsubscript{ven} (mg/L)} &
\textbf{Final age (yr)} \\
\midrule
No breastfeeding & -- & -- & -- & -- \\
6 months BF      & -- & -- & -- & -- \\
1 year BF        & -- & -- & -- & -- \\
3 years BF       & -- & -- & -- & -- \\
\bottomrule
\end{tabular}
\end{table}

% NOTE: Add a paragraph interpreting results:
%  - How does breastfeeding duration affect peak PFAS concentration?
%  - Compare with published epidemiological data (Grandjean 2012, Lau 2012)
%  - Discuss the in-utero contribution to initial body burden at birth

\subsection{FAIR Compliance Evaluation}

% NOTE: This is a key section for the special issue. Add a FAIR checklist
% or evaluation table with columns: FAIR Principle | Criterion | How
% PBKFAIRModel satisfies it. Example rows:
%  - Findable: persistent identifier for SBML model (GitHub DOI via Zenodo)
%  - Accessible: open-source code, web UI without installation
%  - Interoperable: SBML Level 3 V2 format, REST API
%  - Reusable: modular Python package, documented parameters, open license
% Reference the RDA FAIR Maturity Indicators if applicable.

The SBML encoding satisfies the Interoperability and Reusability FAIR
principles by producing a machine-readable, community-standard model
representation compatible with tools such as COPASI and Jarnac. Web-
based deployment via FAIRDatabase satisfies Accessibility by removing
the need for local installation. Persistent storage of simulation
runs with provenance metadata contributes to Findability and
Reusability~\cite{wilkinson2016}.

\subsection{Comparison with Existing Platforms}

% NOTE: This subsection is appropriate for the main text. Expand to
% include a comparison TABLE with columns:
% Platform | Model format | Web interface | Persistent storage |
% FAIR compliance | PFAS support | Life-stage coverage
% Include httk, IRRA, IndusChemFate, and PBPKdb as comparators.

Several existing platforms provide access to toxicokinetic models.
The EPA's \texttt{httk} R package~\cite{pearce2017} provides rapid
tiered assessment for thousands of chemicals but lacks a web-based
interface and persistent storage. The IRRA framework~\cite{mACQUIRE}
offers interactive risk assessment but is limited to reproductive
toxicity endpoints. PBKFAIRModel within FAIRDatabase uniquely combines
SBML-encoded model execution, persistent artifact storage, and
collaborative role-based access control within a single open-source
platform.

\subsection{Limitations and Future Work}

% NOTE: Keep this subsection. Strengthen it by discussing:
%  1. Lack of uncertainty quantification (Bayesian calibration -- cite
%     Bois et al. MCSim for Bayesian PBK context)
%  2. Validation against biomonitoring data from Faroe Islands cohort
%     (Grandjean 2012) as a future step
%  3. Extension to chemical mixtures (aligns with special issue scope)
%  4. Integration with microbiome data in FAIRDatabase (Theme 4: NAMs)

The current implementation has several limitations. Synchronous ODE
execution requires 30--90 seconds per scenario; future work will
implement asynchronous job queuing (e.g., Celery) for production
scalability. The model is currently parameterized for PFOA and PFOS;
expansion to the full PFAS chemical space requires systematic
parameterization from biomonitoring data and high-throughput
predictions. Uncertainty quantification through Bayesian parameter
calibration would substantially enhance the model's utility for
regulatory decision support. Future integration with microbiome and
multi-omics datasets within FAIRDatabase would enable exploration of
host-microbiome interactions in PFAS toxicokinetics, aligning with
New Approach Methodology (NAM) frameworks.

% ====================== CONCLUSION =========================
\section{Conclusion}\label{sec5}

% NOTE: Keep conclusion brief (1 paragraph, ~100 words). Current draft
% is appropriate in length. Ensure the final sentence frames the work
% explicitly in the context of next-generation risk assessment to match
% the special issue title.

We presented PBKFAIRModel, a fully FAIR-compliant physiologically-based
kinetic modelling framework integrated into FAIRDatabase for
computational toxicology. The SBML-encoded lifetime PBK model, deployed
via a Python-native simulation engine and web-based interface, enables
accessible quantification of PFAS toxicokinetics across critical
life stages including pregnancy, lactation, and early childhood.
Role-based access control and persistent storage of simulation runs
support collaborative, reproducible research workflows. By combining
model execution, data management, and FAIR-compliant provenance
tracking within a unified open-source platform, PBKFAIRModel provides
a concrete step toward the next-generation risk assessment infrastructure
called for by ongoing global initiatives including PARC and OECD
PBK model development efforts.

% ====================== DECLARATIONS =======================
\section*{Acknowledgements}
The authors thank the reviewers for their constructive feedback. This
work was conducted within the MetaHealth Project.

% NOTE: Add a Declaration of Generative AI Use section if any AI tools
% were used in manuscript preparation (REQUIRED by Computational Toxicology).
% Example:
% \section*{Declaration of generative AI and AI-assisted technologies
%   in the manuscript preparation process}
% During the preparation of this work the authors used [TOOL] in order
% to [REASON]. After using this tool, the authors reviewed and edited
% the content as needed and take full responsibility for the content.

\section*{Declaration of competing interests}
The authors declare that they have no conflict of interest.

\section*{Funding}
% NOTE: Use the standard Elsevier funding format below.
This work was supported by the Dutch Research Agenda research program
(NWO, NWA-ORC 020/21) -- MetaHealth Project [grant number
NWA.1389.20.080].

\section*{CRediT Author Statement}
\textbf{T.-K. Ly}: Conceptualization, Methodology, Software,
Validation, Writing -- Original Draft.
\textbf{S. Wiedeman}: Methodology, Software, Validation.
\textbf{R. V. Bumbuc}: Conceptualization, Resources,
Writing -- Review \& Editing, Supervision.
\textbf{V. M. Sheraton}: Conceptualization, Writing -- Review \&
Editing, Supervision.

% ====================== DATA STATEMENT ====================
% NOTE: A Data Statement IS REQUIRED by Computational Toxicology.
% State where the model code and simulation data are available.
% Example:
\section*{Data availability}
The SBML model file (\texttt{lifetime\_pbpk.xml}), Python simulation
runner, and FAIRDatabase integration code are available on GitHub at
[URL to be inserted]. A persistent DOI will be registered via Zenodo
upon acceptance. % NOTE: Register a Zenodo DOI for the model BEFORE
                 % submission -- this directly satisfies FAIR Findability
                 % and is expected by the special issue editors.

% ============================================================
% REFERENCES
% NOTE: Computational Toxicology uses NUMBERED citations in square
% brackets [1,2] style (NOT author-year). The current manuscript uses
% \cite{} commands -- these will compile correctly with natbib numbers
% style as set above. Ensure bibliography style is set to:
%   \bibliographystyle{elsarticle-num}
% ============================================================
\section*{References}
\bibliographystyle{elsarticle-num}
\bibliography{references}

\end{linenumbers}
\end{document}
